\documentclass[10pt]{article}
\usepackage[T1]{fontenc}
\usepackage[utf8]{inputenc}
\usepackage[french]{babel}
\usepackage{amsmath,amssymb,mathrsfs,mathpazo}
\usepackage{geometry}
\geometry{a4paper,margin=2cm}
\usepackage{graphicx,array,tabularx,tabularray,colortbl,multirow,multicol,booktabs,enumitem,pifont,ulem,fancybox,diagbox,cancel}
\usepackage{pstricks,pstricks-add,pst-plot,pst-tree,pst-node,pst-text,pst-func,pst-math,pst-bspline,pst-3dplot,pst-eucl,pst-solides3d}
\let\pscalculate\undefined
\usepackage{pst-bezier}
\usepackage{tikz,pgfplots,tkz-tab}
\usetikzlibrary{arrows,arrows.meta,patterns,calc,positioning,trees,decorations.markings,shapes.geometric,shapes.misc,shapes.arrows,intersections,decorations.pathmorphing,decorations.pathreplacing,backgrounds,shadows,babel}
\pgfplotsset{compat=1.18}
\usepgfplotslibrary{fillbetween}
\usepackage[np]{numprint}
\usepackage{listings,makecell,pas-tableur,twoopt}
\let\barre\undefined
\usepackage[pstricks]{tablvar}
\definecolor{darkgreen}{RGB}{0,100,0}
\providecommand{\red}{\color{red}}
\providecommand{\blue}{\color{blue}}
\providecommand{\green}{\color{green}}
\providecommand{\black}{\color{black}}
\providecommand{\white}{\color{white}}
\providecommand{\gray}{\color{gray}}
\providecommand{\cyan}{\color{cyan}}
\providecommand{\magenta}{\color{magenta}}
\providecommand{\yellow}{\color{yellow}}
\providecommand{\violet}{\color{violet}}
\providecommand{\orange}{\color{orange}}
\providecommand{\euro}{\texteuro}
\providecommand{\EUR}{\texteuro}
\providecommand{\ieme}{\textsuperscript{e}}
\providecommand{\degres}{\ensuremath{^\circ}}
\providecommand{\degree}{\ensuremath{^\circ}}
\providecommand{\vv}[1]{\overrightarrow{#1}}
\providecommand{\Coord}[1]{#1}
\providecommand{\decosix}{}
\providecommand{\starredbullet}{\ensuremath{\bullet}}
\pagestyle{empty}
\setlength{\parindent}{0pt}
\usepackage{adjustbox}\DeclareUnicodeCharacter{2236}{:}\DeclareUnicodeCharacter{2190}{\ensuremath{\leftarrow}}
\providecommand{\R}{\mathbb{R}}
\providecommand{\N}{\mathbb{N}}
\providecommand{\Z}{\mathbb{Z}}
\providecommand{\Q}{\mathbb{Q}}
\providecommand{\D}{\mathbb{D}}
\providecommand{\C}{\mathbb{C}}
\providecommand{\e}{{\mathrm{e}}}
\providecommand{\vect}[1]{\overrightarrow{#1}}
\providecommand{\vectt}[1]{\overrightarrow{\mathrm{#1}}}
\providecommand{\vvt}[1]{\overrightarrow{\mathrm{#1}}}
\providecommand{\np}[2][]{{#2~#1}}
\providecommand{\vv}[1]{{\overrightarrow{#1}}}
\renewcommand{\d}{{\mathrm{d}}}
\providecommand{\degres}{{^\circ}}
\providecommand{\degre}{{^\circ}}
\providecommand{\blue}{\color{blue}}
\providecommand{\red}{\color{red}}
\providecommand{\enskip}{\quad}
\providecommand{\llbracket}{\lbrack\!\lbrack}
\providecommand{\rrbracket}{\rbrack\!\rbrack}
\providecommand{\euro}{{\text{€}}}
\providecommand{\no}{{\text{n°}}}
\providecommand{\interval}[2]{{[#1;#2]}}
\providecommand{\intervFO}[2]{{[#1;#2[}}
\providecommand{\dotfill}{\ldots}

\providecommand{\ds}{\displaystyle}
\providecommand{\textding}[1]{#1}
\providecommand{\Oij}{\(\mathrm{O};\vec{\imath},\vec{\jmath}\)}
\providecommand{\Oijk}{\(\mathrm{O};\vec{\imath},\vec{\jmath},\vec{k}\)}
\providecommand{\Ouv}{\(\mathrm{O};\vec{u},\vec{v}\)}
\providecommand{\intervFF}[2]{\ensuremath{[#1;#2]}}
\renewcommand{\euro}{\ifmmode\text{\texteuro}\else\texteuro\fi}
\DeclareUnicodeCharacter{00B0}{\ifmmode^{\circ}\else\textdegree\fi}
\DeclareUnicodeCharacter{20AC}{\ifmmode\text{\texteuro}\else\texteuro\fi}
% Énoncés : Éducation nationale. Transcription et corrigés : APMEP, Denis Vergès et contributeurs. Découpage et classement : CoopMaths, Sébastien Lozano, Éric Elter et contributeurs. Adaptation de la présentation : Maths Explorer.
% CC BY-SA : https://coopmaths.fr/www/about/
% https://coopmaths.fr/alea/static/e3c/2020/tex/e3c_2020_00_sujet10_2_cor.tex
\begin{document}
\noindent\textit{APMEP / CoopMaths. CC BY-SA. Présentation adaptée par Maths Explorer.}\par\medskip
\section*{Exercice 2 (5 points)}

\subsection*{1. Justifier que pour tout réel $x$ appartenant à $[0 ; 9]$ : $V(x) = 4x^3 - 84x^2 + 432x$.}
La longueur est égale à $24 - 2x$ et la largeur $18 - 2x$. La hauteur valant $x$, le volume de la boîte sans couvercle est égal à :
\[V(x) = L \times l \times h = (24 - 2x)(18 - 2x) \times x = x(432 - 48x - 36x + 4x^2) = 4x^3 - 84x^2 + 432x	\]
	
\subsection*{2. On note $V'$ la fonction dérivée de $V$ sur $[0 ; 9]$. Donner l’expression de $V'(x)$ en fonction de $x$.}
La fonction polynôme $V$ est dérivable sur $\mathbb{R}$ et donc sur $[0 ; 9]$,
	\[	V'(x) = 12x^2 - 168x + 432 = 12(x^2 - 14x + 36)	\]
	
\subsection*{3. Dresser alors le tableau de variations de $V$ en détaillant la démarche.}
Avec $\Delta = 14^2 - 4 \times 36 = 196 - 144 = 52$, le trinôme a deux racines :
	\[
	x_1 = \dfrac{14 + \sqrt{52}}{2} = 7 + \sqrt{13} \quad \text{et} \quad x_2 = \dfrac{14 - \sqrt{52}}{2} = 7 - \sqrt{13}
	\]
On sait que ce trinôme est positif sauf sur $]7 - \sqrt{13} ; 7 + \sqrt{13}[$. Comme $\sqrt{13} \approx 3,6$, $7 + \sqrt{13} \approx 10,6$. Donc :
	\begin{itemize}
		\item sur $]7 - \sqrt{13} ; 7 + \sqrt{13}[$, $V'(x) < 0$, donc $V$ est décroissante ;
		\item sur $[0 ; 7 - \sqrt{13}]$ et sur $[7 + \sqrt{13} ; 9]$, $V'(x) > 0$, donc $V$ est croissante.
	\end{itemize}
	
	\subsection*{4. Pour quelle(s) valeur(s) de $x$ la contenance de la boîte est-elle maximale ?}
	
	D’après la question précédente, $V(7 - \sqrt{13}) \approx 654,98 \approx 655 \text{ cm}^3$ est la contenance maximale de la boîte.
	
	\subsection*{5. L’industriel peut-il construire ainsi une boîte dont la contenance est supérieure ou égale à 650 cm$^3$ ? Justifier.}
	
	Oui, puisque la capacité maximale est proche de $655$.\\
	 Pour avoir une capacité de $650$ cm$^3$, il faut prendre $x \approx 3,06$ cm.

\end{document}
